% TODO checklist: Inspected files:
% - README.md (terminology throughout)
% - Q&A.md (definitions in answers)
% - docs/general/GLOSSARY.md
% - src/ (code terminology)

\chapter{Glossary of Terms}
\label{app:glossary}

This appendix defines key terms, acronyms, and concepts used throughout this thesis and the Cineca Agentic Platform documentation.

%------------------------------------------------------------------------------
\section{General Terms}
\label{app:glossary:general}

\begin{description}[style=nextline,leftmargin=2cm,labelwidth=1.8cm]

\item[Agent]
An autonomous software entity that can perceive its environment, make decisions, and take actions to achieve goals. In this platform, agents are LLM-powered systems that execute multi-step workflows.

\item[Agentic AI]
AI systems that exhibit autonomous, goal-directed behavior, capable of planning, reasoning, and using tools to accomplish complex tasks without explicit step-by-step instructions.

\item[Control Plane]
The layer responsible for configuration, orchestration, and management of platform components. In this architecture, PostgreSQL serves as the control plane data store.

\item[Data Plane]
The layer handling low-latency, high-throughput data operations. Redis serves this role for caching, queuing, and session state.

\item[Multi-Tenancy]
An architecture where a single application instance serves multiple tenant organizations with logical data isolation.

\item[Orchestration]
The coordination of multiple components, services, or steps to accomplish a complex task. The Agent Orchestration Engine coordinates LLM reasoning, tool invocation, and state management.

\item[Pipeline]
A sequence of processing stages where the output of one stage becomes the input of the next. The NL$\rightarrow$Cypher pipeline translates natural language to graph queries.

\item[Resilience]
The ability of a system to continue operating correctly in the presence of failures, including LLM provider outages, network issues, and resource exhaustion.

\item[Tool]
A discrete, invokable capability that agents can use to interact with external systems, process data, or perform specialized tasks.

\end{description}

%------------------------------------------------------------------------------
\section{Architecture Terms}
\label{app:glossary:architecture}

\begin{description}[style=nextline,leftmargin=2cm,labelwidth=1.8cm]

\item[ADR]
Architecture Decision Record. A document capturing the context, decision, and consequences of an architectural choice.

\item[API Gateway]
A single entry point for all client requests, providing routing, authentication, rate limiting, and request transformation.

\item[Circuit Breaker]
A fault tolerance pattern that prevents cascading failures by temporarily stopping requests to a failing service after a threshold of failures.

\item[CQRS]
Command Query Responsibility Segregation. A pattern separating read and write operations into different models or services.

\item[Event Sourcing]
An architectural pattern where state changes are stored as a sequence of events, enabling full audit trails and state reconstruction.

\item[Microservices]
An architectural style where an application is composed of loosely coupled, independently deployable services.

\item[Repository Pattern]
A data access pattern that abstracts database operations behind a collection-like interface, promoting testability and decoupling.

\item[SSE]
Server-Sent Events. A protocol for servers to push real-time updates to clients over HTTP. Used for job progress streaming.

\end{description}

%------------------------------------------------------------------------------
\section{LLM and AI Terms}
\label{app:glossary:llm}

\begin{description}[style=nextline,leftmargin=2cm,labelwidth=1.8cm]

\item[Chain-of-Thought (CoT)]
A prompting technique where the LLM is encouraged to articulate its reasoning step-by-step before providing a final answer.

\item[Context Window]
The maximum number of tokens an LLM can process in a single request, including both input and output.

\item[Embedding]
A dense vector representation of text that captures semantic meaning, enabling similarity search and clustering.

\item[Fine-tuning]
The process of training a pre-trained model on a specific dataset to adapt it for particular tasks or domains.

\item[Hallucination]
When an LLM generates plausible-sounding but factually incorrect or unsupported information.

\item[Intent Classification]
The process of determining the user's intended action from their natural language input. Used as a security filter.

\item[LLM]
Large Language Model. A neural network trained on vast text corpora, capable of understanding and generating human-like text.

\item[MCP]
Model Context Protocol. A standardized interface for tool discovery, invocation, and result handling in agentic systems.

\item[Ollama]
An open-source platform for running LLMs locally, providing an OpenAI-compatible API.

\item[Prompt Engineering]
The practice of crafting effective prompts to elicit desired behaviors from LLMs.

\item[RAG]
Retrieval-Augmented Generation. A technique combining information retrieval with LLM generation to provide grounded, factual responses.

\item[Temperature]
A hyperparameter controlling the randomness of LLM outputs. Lower values (e.g., 0.0) produce deterministic outputs; higher values (e.g., 1.0) increase creativity.

\item[Token]
The basic unit of text processed by LLMs. Roughly equivalent to 3-4 characters or 0.75 words in English.

\item[Tool Calling]
The capability of LLMs to invoke external functions or APIs based on natural language instructions.

\end{description}

%------------------------------------------------------------------------------
\section{Database Terms}
\label{app:glossary:database}

\begin{description}[style=nextline,leftmargin=2cm,labelwidth=1.8cm]

\item[Alembic]
A database migration framework for SQLAlchemy, enabling version-controlled schema changes.

\item[Bolt Protocol]
A binary protocol used by graph databases (Neo4j, Memgraph) for efficient client-server communication.

\item[Cypher]
A declarative graph query language used by Memgraph and Neo4j, similar to SQL for graph data.

\item[JSONB]
A PostgreSQL data type for storing JSON documents in a binary, indexed format.

\item[Memgraph]
An in-memory graph database using the Cypher query language, optimized for real-time analytics.

\item[ORM]
Object-Relational Mapping. A technique for accessing relational data using object-oriented paradigms. SQLAlchemy is the ORM used in this platform.

\item[PostgreSQL]
An advanced open-source relational database known for robustness, extensibility, and standards compliance.

\item[Redis]
An in-memory data store supporting strings, lists, sets, sorted sets, and hashes. Used for caching, queues, and pub/sub.

\item[SQLAlchemy]
A Python SQL toolkit and ORM providing database abstraction and query generation.

\end{description}

%------------------------------------------------------------------------------
\section{Security Terms}
\label{app:glossary:security}

\begin{description}[style=nextline,leftmargin=2cm,labelwidth=1.8cm]

\item[JWT]
JSON Web Token. A compact, URL-safe token format for representing claims between parties, used for authentication.

\item[JWKS]
JSON Web Key Set. A set of keys containing public keys used to verify JWT signatures.

\item[OAuth 2.0]
An authorization framework enabling third-party applications to access resources on behalf of a user.

\item[OIDC]
OpenID Connect. An identity layer on top of OAuth 2.0, providing authentication and user identity information.

\item[Output Guard]
A security mechanism that validates LLM outputs for harmful content, policy violations, or sensitive data.

\item[PII]
Personally Identifiable Information. Data that can identify an individual, such as names, emails, or phone numbers. The platform scrubs PII from logs.

\item[RBAC]
Role-Based Access Control. An authorization model where permissions are assigned to roles rather than individual users.

\item[Scope]
A permission unit in OAuth/OIDC defining what resources a token can access (e.g., \texttt{graph:read}, \texttt{admin:all}).

\item[Tenant]
An organizational unit in a multi-tenant system, with isolated data and potentially separate configuration.

\end{description}

%------------------------------------------------------------------------------
\section{Observability Terms}
\label{app:glossary:observability}

\begin{description}[style=nextline,leftmargin=2cm,labelwidth=1.8cm]

\item[APM]
Application Performance Monitoring. Tools and practices for measuring and managing application performance.

\item[Correlation ID]
A unique identifier passed through all services handling a request, enabling distributed trace analysis.

\item[Grafana]
An open-source analytics and visualization platform, used to display metrics from Prometheus.

\item[Health Check]
An endpoint that returns the health status of a service, used by orchestrators for liveness and readiness probes.

\item[Metrics]
Numeric measurements collected over time, such as request counts, latencies, and error rates.

\item[OpenTelemetry (OTel)]
A vendor-neutral observability framework providing APIs and tools for tracing, metrics, and logging.

\item[Prometheus]
A time-series database and monitoring system that scrapes metrics from instrumented services.

\item[Span]
A unit of work in a distributed trace, representing an operation with a start time and duration.

\item[Trace]
A distributed trace following a request across multiple services, composed of spans.

\end{description}

%------------------------------------------------------------------------------
\section{Platform-Specific Terms}
\label{app:glossary:platform}

\begin{description}[style=nextline,leftmargin=2cm,labelwidth=1.8cm]

\item[Agent Run]
A single execution of an agent, from receiving a user request to producing a final response.

\item[Agent Session]
A stateful conversation context containing multiple agent runs (steps) with shared history.

\item[Agent Step]
An individual action within an agent session: user message, assistant response, tool invocation, or system event.

\item[Built-in LLM Manifest]
Configuration for locally-managed LLM processes, including model paths, ports, and resource limits.

\item[Job]
A unit of background work with a lifecycle (queued $\rightarrow$ running $\rightarrow$ finished/failed) and associated payload and result.

\item[Model Instance]
A loaded, configured LLM accessible through a provider, with associated capabilities and settings.

\item[NL$\rightarrow$Cypher]
The pipeline translating Natural Language questions into Cypher graph queries for Memgraph.

\item[Provider]
An LLM service configuration (OpenAI, Ollama, etc.) with connection details and credentials.

\item[Provider Pool]
A collection of LLM providers that can be used for failover and load balancing.

\item[Tool Invocation]
A recorded execution of a tool, including input parameters, output, latency, and any errors.

\end{description}

%------------------------------------------------------------------------------
\section{Acronyms}
\label{app:glossary:acronyms}

\begin{table}[ht]
\centering
\caption{Common Acronyms}
\label{tab:acronyms}
\small
\begin{tabularx}{\textwidth}{lX}
    \hline
    \textbf{Acronym} & \textbf{Meaning} \\
    \hline
    ADR & Architecture Decision Record \\
    API & Application Programming Interface \\
    CORS & Cross-Origin Resource Sharing \\
    CRUD & Create, Read, Update, Delete \\
    DTO & Data Transfer Object \\
    HSTS & HTTP Strict Transport Security \\
    HTTP & Hypertext Transfer Protocol \\
    JSON & JavaScript Object Notation \\
    JWT & JSON Web Token \\
    LLM & Large Language Model \\
    MADR & Markdown Architecture Decision Records \\
    MCP & Model Context Protocol \\
    NL & Natural Language \\
    OIDC & OpenID Connect \\
    ORM & Object-Relational Mapping \\
    OTel & OpenTelemetry \\
    PII & Personally Identifiable Information \\
    RBAC & Role-Based Access Control \\
    REST & Representational State Transfer \\
    RPC & Remote Procedure Call \\
    SQL & Structured Query Language \\
    SSE & Server-Sent Events \\
    TLS & Transport Layer Security \\
    TTL & Time To Live \\
    UI & User Interface \\
    UUID & Universally Unique Identifier \\
    \hline
\end{tabularx}
\end{table}

